\Askhsh{37106}{4}
\begin{ylh}{1}
    \item 4.2. Τέμνουσα δύο ευθειών - Ευκλείδειο αίτημα
    \item 5.6. Εφαρμογές στα τρίγωνα.
\end{ylh}
% ===================================================================
%  Αρχείο: 37106.tex (Fragment)
%  Αυτόματη μετατροπή μέσω doc2latex.py
% ===================================================================

ΘΕΜΑ 4

Δίνεται τρίγωνο ΑΒΓ με ΑΒ \textless{} ΒΓ και η διχοτόμος ΒΕ της γωνίας $\widehat{B}$. Αν ΑΖ $\bot$ ΒΕ, όπου Ζ σημείο της ΒΓ και Μ το μέσον της ΑΓ, να αποδείξετε ότι :

\begin{alist}
\item Το τρίγωνο ΑΒΖ είναι ισοσκελές. \monades{7}

\item ΔΜ // ΒΓ και $\Delta Μ = \ \frac{Β\Gamma - ΑΒ}{2}$. \monades{10}

\item $Ε\widehat{\Delta}Μ = \ \frac{\widehat{Β}}{2}$, όπου $\widehat{B}$ η γωνία του τριγώνου ΑΒΓ. \monades{8}

\includesvg[width=5.43886in,height=3.36458in]{/home/spyros/Μαθηματικά/Trapeza_Thematwn/A_Lykeiou/Geometry/extracted/37106/media/image2.svg}
\end{alist}
